% 第 5 章　怎样描述运动
\chapter{怎样描述运动}

\step{反常识开场}

\begin{mainline}
高铁以每小时 300 千米的速度掠过平原。车厢里，一个男孩把一颗苹果竖直向上抛起，苹果划出一道弧线——落回他的手心。苹果没有“呼”地一声飞向车尾，没有砸中后排的乘客，一切就好像列车停在原地。

这不对劲。想想站台上的朋友看见了什么：男孩抛起苹果的这半秒钟里，列车已经向前冲了四十多米。苹果脱手之后悬在空中，手、座椅、地板都在以每秒八十多米的速度飞奔，它凭什么不落后？是谁在空中推着苹果，陪列车一起飞？

同样奇怪的事每天都在车窗上发生。下雨天留意侧窗玻璃：车停着时，雨痕竖直向下；车一启动，雨痕全都变成斜线，车速越快，雨痕越接近水平。雨明明往下落，为什么在你的窗上横着走？

还有一桩更大的怪事。地球带着我们绕太阳狂奔，速度约每秒 30 千米，是高铁的三百多倍；太阳系又拖着地球绕银河系中心转，速度约每秒 220 千米。我们每时每刻都在以炮弹望尘莫及的速度飞驰，可桌上的咖啡纹丝不动，你原地跳起还是落回原地，从小到大没有任何一个早晨，你感到过这股“每秒二百二十千米的风”。

对比一下就更蹊跷：民航客机巡航时每小时 900 千米，你安稳入睡；过山车时速不过一百出头，你却尖叫到失声。可见身体对“多快”毫无兴趣，它尖叫的对象另有其人。那个对象是谁，是本章要揪出来的主角之一。

这三件怪事其实是同一件事。它们背后藏着物理学里最朴素也最深刻的一个发现：\emph{“运动”不是贴在物体身上的标签，而是物体与观察者之间的一种关系}。这一章要做的，是把“怎样描述运动”这件你以为小学就会的事，重新认真做一遍——做到能解释苹果、雨痕和地球为止。
\end{mainline}

\step{先猜一猜}

往下读之前，请先押上你的直觉。猜错比猜对更值钱，因为本章要修理的正是直觉。

第一题。匀速行驶的火车里，你竖直向上跳起，会落在（a）原处；（b）原处偏后一点；（c）原处偏前一点？如果是火车\emph{刹车}的时候跳呢？

第二题。两个同学争论。甲说：“我坐在高铁上一动不动。”乙说：“胡说，你正以每秒八十多米狂奔。”谁对？能不能设计一个实验，在\emph{不朝窗外看}的前提下，让他们俩分出胜负？

第三题。把一个小球竖直向上抛，到达最高点的那一瞬间，它的速度是多少？它在那一瞬间的“加速的劲头”（本章后面会叫它加速度）又是多少——是零，还是不是零？

写下你的三个答案。本章结尾我们会逐一核对。如果你的答案和多数人一样，那么至少有一题，你的直觉已经站在了物理学家花了两千年才爬出去的那个坑里。

\step{动手实验}

\begin{experiment}[用慢动作视频称一称“下落”]
\textbf{材料}：一部带慢动作摄像模式的手机、一把卷尺或直尺、一枚硬币或一颗橡皮、一面墙。

\textbf{步骤一}：在墙上竖直贴一条纸带，用尺每隔 10 厘米画一条横线，标上 0、10、20、30……这就是你的“坐标尺”。

\textbf{步骤二}：让硬币从刻度 0 处由静止下落，用慢动作模式（每秒 240 帧更好）竖直拍摄，保证整面刻度尺都在画面里。

\textbf{步骤三}：逐帧回放。以硬币脱手那一帧为第 0 帧，记下第 5、10、15、20 帧时硬币所在的刻度。你会发现：相等的时间间隔里，硬币走过的距离\emph{越来越长}——它不是匀速的，它在加速。

\textbf{步骤四}：把手机帧率换算成时间（每秒 240 帧，每帧约 \(1/240\) 秒），把“第几帧”换算成“第几秒”，在方格纸上画出硬币的位置—时间图。点连起来不是直线，而是一条越来越陡的弯线。再算一算：如果下落时间加倍，下落的距离大约变成几倍？多数人会量出一个接近“四倍”的结果——时间两倍，距离四倍，这正是“距离正比于时间的平方”。

\textbf{记录}：把图留下，本章后面会用它量出一个真实的大自然常数。如果你在第三步发现硬币“匀速下落”，多半是按到了暂停键——自由下落的东西从来不匀速。
\end{experiment}

还有一个零成本的观察：下次坐匀速行驶的公交车或地铁时，把钥匙竖直向上抛一小段再接住，连抛五次，注意落点相对你的手有没有偏移。然后请司机师傅（在想象里）急刹车时再做一次——当然，这一次请只在脑子里做，或者用荡秋千时减速的那一下代替。把两次结果记下来，结尾时用得上。

\step{旧模型遇到麻烦}

先替日常语言说句公道话。我们平时说“这车开得真快，每小时一百二”，从来不说明“相对谁”。这套说法能活下来，是因为我们全都站在同一个默认的观察台上：地面。凡是不特别声明，速度一律相对于地面，于是大家相安无事。旧模型可以概括为一句话：\emph{物体的位置和速度是它自己的属性，说清楚了就行，不必提观察者}。

但麻烦一个接一个地来了。

麻烦一：并排列车的错觉。你坐在停靠的列车里，旁边轨道上另一列车缓缓启动。有那么一两秒，你完全分不清是它在向前开，还是你在向后退。你低头看地面才能“破案”——可如果这是在两条平行的传送带上，脚下没有地面呢？旧模型说“谁在动”是客观事实，你的眼睛却告诉你：只看对方，根本分不出来。

麻烦二：谁在超速。导航语音说“您已超速”，它测的是你的车相对\emph{地面}的速度。可对一颗掠过头顶的卫星来说，你和警车都在以每秒数百米的速度随地球自转移动；对太阳来说，你们俩都在以每秒 30 千米绕圈。“快”和“慢”离了观察者，连定义都立不住。

麻烦三：斜着走的雨。站台上的保安看雨竖直下落，车里的你看雨斜着扑面而来。同一场雨，两种轨迹，\emph{两种描述都对}。旧模型问“雨到底朝哪个方向落”，自然界拒绝回答——它只回答“相对于谁”。

麻烦四：太阳到底动不动。每天清晨，你看见太阳从地平线上升起——一个完全合格的运动描述，以你为参考系。哥白尼说，同一件事也可以描述为：太阳没动，是你脚下的地面在向下转动，把你转进了阳光里。两种描述预言的日出时间分秒不差。旧模型急着问“到底谁对谁错”，可“到底”二字恰恰是多出来的：运动学只管相对关系，不问谁的运动是“真身”。选哪个参考系是方便与否的问题，不是真假问题——当然，到了讨论“为什么会这样动”的动力学层面，有些参考系会明显比其他参考系讲理得多，这是后话。

麻烦五：一个哲学的窟窿。有人会说：那就规定一个“真正静止”的裁判席，一切速度相对它来报。牛顿当年正是这样设想的，他管它叫绝对空间。问题在于伽利略发现的一条死规矩：在一间关死的、匀速前进的船舱里，你做任何力学实验——抛球、滴水、摆钟——结果都和船停着时一模一样。如果“绝对静止”存在，它理应能在这类实验里露出马脚；可它从不露面。一个永远测不到的“裁判席”，对物理学来说等于不存在。

五条麻烦指向同一个诊断：旧模型错在把“运动”当成了物体自带的标签。新的出路不是去找更权威的裁判，而是反过来承认——\emph{位置和速度从定义起就必须带上观察者}。这需要发明几个今天看来天经地义、当年却石破天惊的概念。
\step{发明新概念}

\begin{mainline}
\textbf{第一个新概念：参考系。}要描述运动，先指定三样东西：一个观察者站在哪儿（原点），用什么尺子量（坐标轴），用什么钟计时。这一套合起来叫\emph{参考系}。从此“物体在哪里”这句话有了严格含义：它是相对某个参考系读出的一组数。位置像什么？像地图上的经纬度——必须先约定本初子午线和赤道，一个地点才有坐标。它不像什么？它不像人的身高，身高不随你换地图而变化，位置的读数却会随参考系彻底改写。

\textbf{第二个新概念：位移与路程不是一回事。}路程是你实际走过的总长度，位移是从起点指向终点的有向线段（这是第 4 章讲的向量）。你绕操场跑一整圈，路程 400 米，位移是零。打车软件按路程收费，因为它关心轮胎滚了多少；而导弹的制导关心位移，因为它只问“离目标还差多远、朝哪个方向”。混淆这两个量是运动学里最低级也最常见的错误。

其实“位置是一组读数”这件事，你每天都在用。电梯里的楼层显示就是一个一维坐标系：原点在地面层，向上为正，地下一层是负数——坐标可以取负，而“负一层”谁也不觉得奇怪。导航软件把整座城市铺成一张二维坐标网，你的定位就是网上的一个点。坐标系不是数学家的发明，它只是把我们本来就有的报位置的习惯，擦得干净了而已。

\textbf{第三个新概念：速度是位置变化的快慢，而且有方向。}同样的时间谁的位置变得多，谁就快。关键的一步是区分两种速度：全程平均下来的叫\emph{平均速度}，某一瞬间的叫\emph{瞬时速度}——汽车仪表盘上跳动的那个数，就是瞬时速度的大小。博尔特百米跑出 9 秒 58，平均速度约每秒 10.4 米，可他后程瞬时速度接近每秒 12 米，起跑瞬间则是零。一个“平均”掩盖了全部剧情。

导航软件的“预计到达时间”就是平均速度的生意：它用剩余路程除以你近来的平均速度，再按路况修修补补。你一路飙到时速 120，它不会提前多少，因为平均速度被前面堵车的十分钟死死压着；反过来，哪怕你中途停车买杯咖啡，只要全程平均速度没变，预计时间也懒得改。平均速度和瞬时速度各有各的诚实，问题只在于你拿哪一个去回答哪一个问题。

\textbf{第四个新概念：加速度是速度变化的快慢。}速度大小变了（起步、刹车）是加速，方向变了（拐弯）也是加速。公交车起步你向后仰、刹车你向前倾、转弯你被甩向一侧——三次“袭击”来自同一个量。加速度是本章埋得最深的一颗种子：下一章你会发现，它才是力真正管辖的对象，运动本身从来不需要谁来维持。

\textbf{第五个新概念：状态。}把“此刻在哪里”和“此刻速度多大、朝哪”两条信息装在一起，就得到运动物体的\emph{状态}。为什么两条缺一不可？看台球：母球停在同一个位置，轻轻推它和发力击球，球桌上的下一幕天差地别——位置相同，未来不同，差别全在速度里。知道状态，加上规律，就能推算下一刻——这正是全书反复追问的三个问题的第一次亮相：系统此刻是什么状态，什么规律让它变，我们怎样测量它。

\textbf{第六个新概念：伽利略相对性原理。}在所有做匀速直线运动、互不加速的参考系（叫惯性参考系）里，力学实验的结果完全相同。船上的人测出的下落规律、碰撞规律，和岸上人测出的一字不差；没有任何实验能告诉你“你的船到底动不动”。这句话反过来说更锋利：\emph{匀速运动是感觉不到的，能被察觉的只有运动的改变}。
\end{mainline}

\begin{center}
\begin{tikzpicture}[scale=0.92]
  % 地面参考系
  \draw[thick] (-5.4,0) -- (5.4,0);
  \fill[pattern=north east lines] (-5.4,-0.26) rectangle (5.4,0);
  \node[below=8pt] at (-4.6,0) {地面参考系};
  % 火车
  \draw[thick,fill=gray!12] (-1.2,0) rectangle (3.6,1.3);
  \draw[thick,fill=white] (-0.9,0.35) circle (0.28);
  \draw[thick,fill=white] (2.9,0.35) circle (0.28);
  \node at (1.2,0.85) {匀速火车};
  \draw[-{Stealth[length=3mm]},very thick,blue!70!black] (3.7,1.6) -- (5.3,1.6) node[right=1pt] {\(v\)};
  % 人与苹果
  \draw[thick,fill=orange!20] (0.9,0.55) circle (0.16);
  \node[below=1pt] at (0.9,0.55) {\scriptsize 人};
  \draw[thick,fill=red!40] (1.9,0.75) circle (0.09);
  % 车上看：竖直
  \draw[dashed,red!70!black] (1.9,0.55) -- (1.9,1.85);
  \node[red!70!black,above=0pt] at (1.9,1.85) {\scriptsize 车上看：直上直下};
  % 地面看：抛物线
  \draw[dashed,blue!70!black] (1.9,0.75) parabola bend (2.7,1.7) (3.5,0.75);
  \node[blue!70!black,below=2pt] at (3.9,0.6) {\scriptsize 地面看：抛物线};
\end{tikzpicture}
\end{center}

这张图已经预告了本章结尾：同一颗苹果，在车上的人看来是直上直下，在站台的人看来是一条抛物线。两种描述都对，而且由同一条规律推出——这正是“参考系”这个概念的威力：它不允许两种描述打架。

\step{一句话模型}

\begin{oneline}
运动没有绝对标签：先报出你的参考系，位置和速度才有意义；而所有匀速前进的观察者彼此平等——谁也别想用任何力学实验证明自己“真的在动”。
\end{oneline}

背下这一句，火车上的苹果、斜落的雨、狂奔的地球就都已经解释了一半；剩下的一半，交给数学。

\step{数学翻译器}

\begin{mainline}
先翻译速度。位置记作 \(x\)，时间记作 \(t\)。一段位移 \(\Delta x\) 花掉时间 \(\Delta t\)，平均速度就是
\[
  \bar v=\frac{\Delta x}{\Delta t}.
\]
读法：\(\Delta\) 念“变化量”，整个式子读作“位置的变化量除以花掉的时间”。做例行检查：若物体不动，\(\Delta x=0\)，则 \(\bar v=0\)，成立；若 \(\Delta x\) 取负（朝反方向走），\(\bar v\) 为负，方向信息自动保留；时间越短，同样的位移给出的速度越大，符合直觉。把 \(\Delta t\) 越缩越小，缩到“一个瞬间”，平均速度就变成了瞬时速度——在位置—时间图上（第 3 章的工具），平均速度是割线的斜率，瞬时速度是切线的斜率。\textbf{斜率就是“变化有多快”}，这张图值得再看一眼：
\end{mainline}

\begin{center}
\begin{tikzpicture}[scale=1.05]
  \draw[-{Stealth[length=2.5mm]}] (0,0) -- (6.4,0) node[below left=2pt] {\(t\)};
  \draw[-{Stealth[length=2.5mm]}] (0,0) -- (0,3.6) node[below left=2pt] {\(x\)};
  % 匀加速曲线
  \draw[very thick,blue!70!black] (0,0) parabola (5.4,3.1);
  \node[blue!70!black,right=2pt] at (5.0,2.55) {匀加速};
  % 匀速直线
  \draw[very thick,red!70!black] (0,0) -- (5.9,1.55);
  \node[red!70!black,right=2pt] at (5.9,1.55) {匀速};
  % 切线示意
  \draw[dashed,gray] (2.4,0.47) -- (4.6,2.6);
  \node[gray,below=1pt] at (4.0,2.2) {\scriptsize 某瞬间的切线斜率 \(=\) 瞬时速度};
  \fill (3.5,1.19) circle (0.06);
\end{tikzpicture}
\end{center}

\begin{mainline}
接着翻译加速度：
\[
  a=\frac{\Delta v}{\Delta t},
\]
即“速度的变化量除以花掉的时间”。检查：匀速运动 \(\Delta v=0\)，故 \(a=0\)——匀速再快也没有加速度，坐在每小时 300 千米的匀速高铁里，你的加速度是零，这就是你感觉不到它的原因；倒车时速度为负且越来越负，\(\Delta v<0\)，加速度与“前进方向”相反，式子照样工作。

匀加速（\(a\) 不变）是最简单也最有用的一类。设初速度为 \(v_0\)，三条公式管全部：
\[
  v=v_0+at,\qquad
  x=x_0+v_0 t+\tfrac12 a t^2,\qquad
  v^2=v_0^2+2a\,(x-x_0).
\]
逐条过检查。令 \(a=0\)：第一条退化为 \(v=v_0\)（匀速），第二条退化为匀速直线运动的 \(x=x_0+v_0 t\)，第三条变成 \(v=v_0\)——三条公式在“没有加速度”时彼此自洽。令 \(t=0\)：第一条给出 \(v=v_0\)，第二条给出 \(x=x_0\)，都回到起点，成立。令 \(a\) 反向（刹车）：第一条预言速度不断减小直到为零再变负——汽车刹停后若不挂倒挡不会真变负，这提醒我们公式只管“刹车那段过程”，不管刹停之后。方向检查全部通过。

\textbf{带真实数据的估算：刹车距离。}干燥路面上，小汽车急刹车的加速度大小约为 \(5\) 米每二次方秒（约半个 \(g\)）。车速 \(50\) 千米每小时，即约每秒 \(13.9\) 米。用第三条公式，令末速度 \(v=0\)：
\[
  0=v_0^2-2a\,s\quad\Rightarrow\quad s=\frac{v_0^2}{2a}=\frac{13.9^2}{2\times5}\approx 19\ \text{米}.
\]
车速翻倍到 \(100\) 千米每小时呢？速度在公式里是\emph{平方}出现的，距离变四倍：约 \(77\) 米——还没算司机约一秒的反应时间里车又裸奔了 \(28\) 米。这就是为什么“车速快一倍”不是“危险大一倍”，而是危险大约大四倍。高速公路上“保持车距”的牌子，背面写的就是这条公式。

反向用一次同一组公式，看看工程里怎样\emph{故意}把加速度做小。高铁从静止加速到每小时 300 千米（约每秒 83 米），司机不会一脚到底，而是让列车用大约 5 分钟从容提速。平均加速度约为
\[
  a=\frac{\Delta v}{\Delta t}=\frac{83\ \text{米/秒}}{300\ \text{秒}}\approx 0.28\ \text{米/秒}^2,
\]
不到 \(g\) 的百分之三。这正是你在车厢里能安稳放一杯咖啡的原因：乘客的舒适度指标，本质上是给加速度设上限，而不是给速度设上限。过山车则反过来经营加速度——它的刺激感几乎全部来自加速度的大小和方向的剧烈变化，而不是来自速度本身。游乐园和高铁，一个把 \(a\) 往大里做，一个往小里做，读的是同一条公式。

你的手机里就住着一只加速度的测量员。那颗指甲盖大小的加速度计芯片，里面悬着一小块质量，手机加速时它相对外壳轻轻偏移，电路把偏移量翻译成三个方向的加速度读数——手机的横竖屏切换、计步、游戏里的重力感应，全靠它。有趣的是，它安静地躺在桌面上时，读数不是零，而是方向朝上的 \(9.8\)：它测到的其实不是“运动的改变”，而是“支持它的东西怎样推它”。这桩看似矛盾的小事，要到第 6、7 章讲清力之后才能彻底了结，第 41 章还会让它在一个更大的舞台上重新登场。

自由落体是匀加速的特例：在地面附近、忽略空气阻力时，一切物体下落的加速度都是同一个常数 \(g\approx 9.8\) 米每二次方秒，方向竖直向下。回到本章实验：你的硬币从静止下落，若用 \(x=\tfrac12 g t^2\) 拟合你画的那条弯线，就能亲手量出 \(g\)。比如下落 1 米，预言用时 \(t=\sqrt{2\times1/9.8}\approx 0.45\) 秒——去看看你的慢动作视频，第 100 帧上下（每秒 240 帧时）硬币是否正好落到 1 米线。
\end{mainline}

\begin{mathbridge}
那三条公式从哪里来？第二条里那个神秘的 \(\tfrac12\) 最值得交代。初速度 \(v_0\)、匀加速到 \(v\)，由于速度\emph{均匀}增长，全程的平均速度恰是首尾的平均：\(\bar v=(v_0+v)/2\)。于是位移
\[
  \Delta x=\bar v\,t=\frac{v_0+(v_0+at)}{2}\,t=v_0 t+\tfrac12 a t^2.
\]
那个 \(\tfrac12\) 不是凑出来的，它是“均匀增长的平均值等于首尾平均”的几何事实——在速度—时间图上，位移是图线与时间轴围成的面积，一个矩形加一个三角形，三角形的面积自然带着因子 \(\tfrac12\)。第三条公式则由前两式消去 \(t\) 得到。想再看深一层：把 \(\Delta t\) 趋于零的极限写成导数，\(v=\dfrac{dx}{dt}\)，\(a=\dfrac{dv}{dt}\)，位置、速度、加速度排成一条“求导链”，下一章会用到这条链。

抛体运动也在这里一次解决。把初速度沿水平和竖直方向分解（第 4 章的向量分解）：水平方向没有任何推拉（忽略空气阻力），速度保持 \(v_{0x}\) 不变；竖直方向只受 \(g\)，做匀加速。两个方向各自独立、同时进行：
\[
  x=v_{0x}\,t,\qquad y=v_{0y}\,t-\tfrac12 g t^2.
\]
消去 \(t\) 得到 \(y\) 关于 \(x\) 的二次函数——抛物线。检查：令抛射角为 \(90^\circ\)（竖直上抛），\(v_{0x}=0\)，剩下 \(y=v_{0y}t-\tfrac12gt^2\)，退化为竖直匀加速；令 \(g\to0\)（想象关掉重力的行星），\(y=v_{0y}t\)、\(x=v_{0x}t\)，合成一条直线——没有重力，抛体走直线，与第一定律式的一致。最高点的速度不是零：竖直分量是零，水平分量还是 \(v_{0x}\)。

生活中到处是这条抛物线。篮球离手之后，水平方向匀速、竖直方向先减后增，球心描出的正是它；广场喷泉里每一股水柱、消防水枪喷出的水龙，都是一串串水滴接力画出的抛物线。为什么投篮要追求较高的出手弧度？抛物线越陡地落下，球筐在球看来“张开”的有效面积越大——同样的出手速度误差，陡弧度的球更宽容。打球的人不懂公式，手感却早已把这条公式算了成千上万遍。
\end{mathbridge}

\begin{center}
\begin{tikzpicture}[scale=1.05]
  \draw[-{Stealth[length=2.5mm]}] (0,0) -- (7.0,0) node[below left=2pt] {\(x\)};
  \draw[-{Stealth[length=2.5mm]}] (0,0) -- (0,3.2) node[below left=2pt] {\(y\)};
  \draw[very thick,domain=0:6.4,smooth,variable=\t,blue!70!black] plot ({\t},{1.5*\t-0.195*\t*\t});
  % 起点速度分解
  \draw[-{Stealth[length=2.5mm]},very thick,red!70!black] (0,0) -- (1.7,1.28) node[above right=0pt] {\(v_0\)};
  \draw[-{Stealth[length=2.5mm]},thick,red!70!black] (0,0) -- (1.7,0) node[below=1pt] {\(v_{0x}\)};
  \draw[-{Stealth[length=2.5mm]},thick,red!70!black] (0,0) -- (0,1.28) node[left=1pt] {\(v_{0y}\)};
  \draw[dashed,gray] (1.7,0) -- (1.7,1.28) -- (0,1.28);
  % 最高点
  \fill (3.85,2.88) circle (0.07);
  \draw[-{Stealth[length=2.5mm]},thick,green!50!black] (3.85,2.88) -- (5.35,2.88) node[right=1pt] {\scriptsize 最高点只剩水平速度};
  \draw[dashed,gray] (3.85,0) -- (3.85,2.88);
\end{tikzpicture}
\end{center}

\begin{mainline}
最后翻译参考系之间的对话。设火车相对地面以速度 \(u\) 前进，你在车上以相对火车 \(v'\) 的速度朝车头走，那么地面测得你的速度为
\[
  v=u+v'.
\]
速度按参考系简单相加，这叫伽利略速度合成。检查：你朝车尾走，\(v'\) 取负，\(v<u\)，成立；若你恰好以 \(v'=-u\) 朝车尾走，\(v=0\)——地面上的朋友会看到你悬在半空中向后退出车站，而你自己感觉只是在车厢里散步。

自动扶梯是这条式子最便宜的实验室。扶梯以每秒约 0.5 米的速度上行，你站着不动，地面测你就是每秒 0.5 米；你以相对扶梯每秒 1 米向上走，地面测你就是每秒 1.5 米；你转身以同样的步速向下走，若步速恰好等于梯速，你就“定”在了原地——旁边商店的服务员看着你原地踏步，你自己却觉得走得挺努力。同一双腿，两个参考系，两种完全不同的“运动故事”，谁也不撒谎。车窗上的雨痕也是同一条式子：雨相对地面竖直向下，车相对地面水平向前，雨相对\emph{车}的速度是两者按向量相减——于是竖直的雨在车窗坐标里带上了水平分量，划出斜线，车速越快斜得越平。开场怪事之二，解决。
\end{mainline}

\begin{challenge}
把“换参考系”写成完整的坐标翻译。地面系的坐标为 \((x,t)\)，火车系以速度 \(u\) 沿 \(x\) 方向行驶，则伽利略变换为
\[
  x'=x-ut,\qquad t'=t.
\]
对 \(t\) 求两次导数：\(v'=v-u\)，\(a'=a\)。注意最后这一行——\emph{速度因参考系而变，加速度不随惯性参考系改变}。这就是为什么“哪些东西是相对的、哪些东西是大家一致的”不能凭感觉乱说：位置、速度是相对的，加速度和时间在牛顿的世界里是大家一致的。第 6 章的力正是和加速度绑定，所以牛顿定律在所有惯性系里长得一模一样——伽利略相对性原理的数学内核就在 \(a'=a\) 这一行里。

再往前看一眼边界。这条变换藏着一个无人察觉的假设：\(t'=t\)，即“两个参考系共用同一座钟”。它在日常生活里精确到挑不出错，但当速度接近光速时，实验会强迫我们修改它——那是第 38 章的故事。本章的公式全是“日常速度极限”下的近似，近似得好到整个工业文明都建立其上。

顺带一提空气阻力。真实的下落满足的规律大致是 \(a=g-kv/m\)：速度越大，阻力越大，加速度越小，最后速度趋向一个极限 \(v_{\text{终端}}=mg/k\)，雨滴正是以每秒几米的终端速度匀速落地的。你在实验里拍硬币时它还很接近纯自由落体；换成一片羽毛，这条修正立刻从“细节”变成“主角”。模型就是这样被一点点逼近真相的。
\end{challenge}

\step{模型的边界}

本章模型的适用条件，值得像念说明书一样念一遍。

第一，我们把物体当成了\emph{质点}：一个有质量、没形状的点。描述列车从北京到上海的运行，列车的 200 米身长无关紧要，质点模型好得很；描述体操运动员的空翻，质点模型立刻报废——身体各部位的运动不一样，那时需要第 11 章的转动语言。模型没有对错，只有合用不合用。

第二，匀加速公式要求 \(a\) 是常数。自由落体在地面附近近似如此，但坐过山车时加速度瞬息万变，公式 \(v=v_0+at\) 一步都不能直接用，得把过程切成无数小段（数学桥层里那条求导链反过来走）或者改用第 9 章的能量方法。

第三，忽略空气阻力是一条假设，不是事实。铅球和硬币的下落符合得很好；羽毛球、雨滴、降落伞完全不吃这一套。伽利略说“一切物体下落一样快”，成立的地方是真空中——阿波罗 15 号的宇航员在月球表面同时松开锤子和羽毛，它们并肩落地，这是这条定律最著名的一次公演。

第四，地面参考系只是\emph{近似}的惯性系。地球在自转，严格说地面上的我们每时每刻都在拐弯，加速度并不为零，只是它造成的偏差在日常生活里小得可以忽略——要小到多小的实验才露馅，是后面的章节再算的事。

第五，最容易被忘掉的一条：本章所有速度都默认远小于光速。这不是本章的疏忽，而是整个牛顿力学的身份证。

第六，关于“瞬时”二字本身。任何真实的测量都要花时间：雷达测速是量一小段时间里的位移，手机加速度计是对质量块的偏移做平均。所以“某一瞬间的速度”永远不是直接测出来的，而是把测量的时间窗越缩越短时，读数趋向的那个极限值。瞬时值是模型的语言，平均值才是测量的语言——两者在窗足够小时几乎重合，日常尽可混用，但心里要知道谁是理想化。

\step{回头解释开场现象}

现在结案。

\textbf{苹果为什么不落后。}男孩抛出苹果之前，苹果已经在以每秒八十多米的速度随列车前进。抛的动作只给了它竖直方向的速度，水平方向没有任何东西推它或拖它——于是按照第 6 章将严格表述的惯性原理，它的水平速度原封不动地保持八十多米每秒。人、苹果、车厢在水平方向上齐头并进，竖直方向苹果只做一次普通的上抛。车上的人看：直上直下；站台的人看：一条每秒水平移动八十多米的抛物线。两种描述通过伽利略速度合成互相翻译，严丝合缝。开头猜一猜的第一题：匀速时落回原处；刹车时，车厢减速而你的水平速度暂时不减，于是落在原处偏\emph{前}。

\textbf{雨为什么横着走。}车窗上雨痕的方向，是“雨相对地面的速度”减去“车相对地面的速度”合成出来的，车速越大，水平分量越大，雨痕越平。

\textbf{地球的风为什么不存在。}我们、大气、咖啡杯全部随地球一起做几乎匀速的运动。根据伽利略相对性原理，匀速运动不产生任何能被本地实验察觉的效应——能感觉到的只有运动的\emph{改变}。地球公转的轨道其实是弯的，我们并非严格匀速，但那点“拐弯的劲头”摊在一年一圈的尺度上，微弱到日常完全无感。“每秒 220 千米”本身，什么感觉都不会制造。

开场问题的完整答案也一并到手了：在匀速行驶的火车里竖直跳起不落到后面，因为“跳”只改变竖直方向的运动，而水平方向你和火车本来就是一体的，并且没有任何机制会让你们分开。真正需要解释的不是“你为什么不落后”，而是“你为什么会以为应该落后”——那是旧模型在你脑子里留下的错觉。

顺手把开头赌桌上的另外两题也结了。第二题：甲和乙\emph{都对}——甲的报告以车厢为参考系，乙的报告以地面为参考系，两份报告之间可以用伽利略速度合成互相翻译，谁也不拥有“真正的那份”。而在不看窗外、只听车厢内部动静的前提下，任何实验都分不出胜负，这正是伽利略相对性原理的判决。第三题先欠着，它藏在下面的挑战 1 里——如果你此刻已经能脱口说出答案，说明本章的新概念真的长在了你身上。

\step{脑洞实验与挑战题}

\begin{thought}[伽利略的无窗船舱]
想象你待在一艘大船的底舱，门窗全部封死，船体隔音避震完美。船在平静的海面上匀速直线行驶。给你全套实验设备：小球、单摆、水滴、弹簧秤、天平。能不能设计一个力学实验，测出船的速度？

伽利略的答案是：不能。舱里的一切实验——球下落的时间、摆的周期、水滴的形状——都与船静止时一模一样。你可以试着把舱想象成以每秒 30 千米绕太阳公转的地球：我们就住在这样一间无窗舱里，而直到望远镜和星空给了我们“窗外”，人类才确信地球在动。

把尺度再推极端一点。想象一艘未来飞船，以每秒 220 千米的速度匀速巡航——和太阳系绕银河的速度相当。舱内的你喝咖啡、抛小球、称体重，一切都和停泊在港口时没有任何区别。速度这个数字再大，只要它不变，你的每一个细胞、每一件仪器都无从知晓。大自然似乎在用最大的嗓门重复同一句话：它只管“变化”，不管“匀速”。

再进一步：如果这个原理是对的，那么“光”这种运动遵守它吗？在舱里测光速，会得到和岸上一样的数吗？这个看似无害的问题，三百年后把“\(t'=t\)”炸得粉碎，逼出了狭义相对论（第 38 章）。本章的模型站在那座革命的门口。
\end{thought}

\begin{puzzles}
\item 竖直上抛的小球到达最高点的那一瞬间，速度为零。它的加速度也是零吗？如果那一瞬间加速度真是零，接下来小球会怎样？\\
\emph{思路提示}：加速度管的是“速度正在变多快”，不是“速度此刻是多少”。做一个思想判决：若最高点加速度为零，按本章公式小球下一瞬间的速度该是多少——它将永远悬在天上。现实不答应。
\item 甲、乙两车在同一直路上行驶，把它们的速度—时间图画在同一张图上，两条曲线在某一时刻相交。这是否意味着两车在那一刻相遇（位置相同）？\\
\emph{思路提示}：速度—时间图记录的是“跑多快”，不是“在哪里”。交点只说明那一刻两车速度相同。位置的信息藏在图线与时间轴围成的\emph{面积}里，还要加上各自的出发点。
\item 一列火车以每秒 70 米匀速前进，车长 200 米。你从车尾以相对车厢每秒 2 米的速度走向车头。当你到达车头时，站台上的观察者测得你移动了多远？\\
\emph{思路提示}：先算你在车上走了多久，再用伽利略速度合成求你相对地面的速度。也可以反过来：这段时间车头前进了多少，你正好比车头少走一个车长。两条路必须给出同一个数——如果不一样，说明哪里算错了。
\item 两列高铁相向而行，时速都是 300 千米。错车的一瞬间，你在其中一列里看另一列，它“嗖”地一下掠过，相对速度约每小时 600 千米——接近民航客机。可为什么两列车里的乘客都安然无恙，连咖啡都不晃一下？再想一想：如果其中一列迎面撞上一堵“以时速 300 千米开过来的墙”，结果和撞上一堵静止的墙相比，哪个更惨？\\
\emph{思路提示}：伤害来自运动状态被迫改变的剧烈程度，而不是相对某个路过物体的速度读数。先选定车上乘客为研究对象，问他的速度在一瞬间被改变了多少。至于那堵“开过来的墙”——是谁在给墙提供时速 300 千米的动力？如果它真能自己开过来，那它和一列火车有什么分别？
\end{puzzles}
